% !TeX root = tesis.tex

\subsection{Funciones de hash criptográficas}
\label{subsec:apendice-a-hashes-funciones-de-hash-criptograficas}

Las funciones de hash son una primitiva criptográfica sumamente popular. Ejemplos conocidos son \MDFive, \SHAOne, \SHATwoFiveSix, \BLAKETwo, \BLAKEThree, \RIPEMD, \PoseidonHash, entre muchas otras. Estas serán utilizadas frecuentemente en el desarrollo del proyecto, y permiten construir primitivas criptográficas más avanzadas.

\begin{definition}[Función de hash criptográfica]
\label{def:apendice-a-hashes-funcion-de-hash-criptografica}
	Sea $n \in \mathbb{N}$. Se dice que $\Hash: \Strings{*} \to \Strings{n}$ es una función de hash criptográfica si satisface las siguientes propiedades:
	\begin{itemize}[noitemsep]
		\item Pueden ser implementadas determinísticamente en forma eficiente.
		\item Resistencia a preimagen: Dado $y \in \Strings{n}$, es difícil encontrar $x \in \Strings{*}$ tal que $\Hash(x) = y$.
		\item Resistencia a segunda preimagen: Dados $x \in \Strings{*}, \Hash(x) \in \Strings{n}$, es difícil encontrar $x' \neq x$ tal que $\Hash(x) = \Hash(x')$.
		\item Resistencia a colisiones: Es difícil encontrar $x \neq x'$ tales que $\Hash(x) = \Hash(x')$.
	\end{itemize}
\end{definition}

La definición anterior es deliberadamente informal, ya que todavía no precisa qué significa \textit{difícil}. En este apéndice alcanza con la intuición de \textit{computacionalmente prohibitivo}; las formulaciones rigurosas de estas propiedades pueden encontrarse en \cite{hac,katz}.

Existe una relación jerárquica entre estas propiedades. En particular, la resistencia a colisiones implica resistencia a segunda preimagen, y esta última implica resistencia a preimagen.

La compresión de un dominio infinitamente grande $\Strings{*}$ en un codominio finito $\Strings{n}$ implica inevitablemente la existencia de colisiones. La seguridad no radica en su inexistencia, sino en la imposibilidad computacional de encontrarlas.

En criptografía, la seguridad de muchas construcciones se analiza en modelos idealizados. En particular, las funciones de hash suelen tratarse no como objetos matemáticos explícitos, sino como primitivas ideales. Por ello, numerosas demostraciones se formulan en modelos como el Random Oracle Model (\ROM) \cite{Bellare1993}. En este marco, se supone que $\Hash: \Strings{*} \to \Strings{n}$ se comporta como una función completamente aleatoria. Bajo esa idealización, cualquier algoritmo genérico para hallar una preimagen tiene complejidad $\BigTheta(2^n)$, mientras que cualquier algoritmo genérico para encontrar colisiones tiene complejidad $\BigTheta(2^{n/2})$ (ver \cite{hac}).

Podemos transformar cualquier estructura matemática o información finita a una tira de bytes, es decir al conjunto $\Strings{*}$, y elegir un $n$ suficientemente grande como para que valga la propiedad de resistencia a colisiones. Esto es lo que se hace en la práctica. Entonces, dada una función de hash $\Hash$ y un conjunto arbitrario $X$ donde sus elementos tienen tamaño finito, podemos hablar de la función de hash asociada $\Hash_X: X \to \Strings{n}$, y por comodidad de notación usar simplemente $\Hash$.

\subsection{Domain separation}
\label{subsec:apendice-a-hashes-domain-separation}

En construcciones criptográficas modernas es habitual reutilizar una misma función de hash en múltiples contextos dentro de un protocolo. Sin embargo, esta reutilización puede introducir interacciones indeseadas entre componentes que, idealmente, deberían comportarse de manera independiente. La técnica de \textit{domain separation} consiste en forzar dicha independencia lógica modificando la entrada de la función de hash mediante etiquetas, prefijos o parámetros públicos.

Es decir, permite que una misma función $\Hash$ actúe como si fueran múltiples funciones independientes. Por ejemplo, en lugar de computar simplemente $\Hash(x)$, se evalúan expresiones del tipo $\Hash(\texttt{commit} \concat x)$ o $\Hash(\texttt{challenge} \concat x)$. De esta forma, aunque los valores de entrada coincidan, los hashes pertenecen a dominios criptográficamente aislados. Esta técnica resulta fundamental a lo largo de esta tesis, por ejemplo, en la heurística de \FiatShamir.

\subsection{Hash hacia curvas elípticas}
\label{subsec:apendice-a-hashes-hash-hacia-curvas-elipticas}

En numerosos protocolos criptográficos modernos resulta necesario mapear datos arbitrarios (es decir, tiras de bits de $\Strings{*}$) a puntos de una curva elíptica. Este problema no es trivial: una función de hash produce cadenas de bits, mientras que los puntos de una curva elíptica deben satisfacer una ecuación algebraica específica. Formalmente, buscamos una función determinística y eficiente

$$\HashToCurve : \Strings{*} \to \E(\Fq)$$

que modele, idealmente, una distribución uniforme sobre un subgrupo criptográficamente relevante de $\E(\Fq)$  \cite{Fouque2012, Wahby2019}. Es decir, que ambas distribuciones (la del hash y la uniforme) sean computacionalmente indistinguibles.

En esta tesis asumiremos la existencia de funciones HashToCurve seguras y eficientes, sin profundizar en sus construcciones explícitas. El estándar moderno que formaliza estas construcciones es RFC 9380 \cite{rfc9380}, que especifica algoritmos concretos. Conceptualmente, dichos algoritmos implementan funciones del tipo

$$\Strings{*} \to \Fq \to \E(\Fq)$$

minimizando sesgos estadísticos \cite{Fouque2012} y evitando ataques estructurales.

En la práctica, estas funciones no se construyen directamente, sino mediante transformaciones intermedias. Típicamente, el proceso involucra:

\begin{enumerate}[noitemsep]
	\item Expandir la entrada mediante una función de hash convencional.
		\item Interpretar el \textit{output} como uno o más elementos del cuerpo base $\Fq$ de la curva.
	\item Aplicar un mapeo algebraico hacia la curva \cite{Wahby2019, cryptoeprint:2013/325}, es decir una función $f: \Fq \to \E(\Fq)$.
	\item (Opcionalmente) Forzar la pertenencia al subgrupo de orden primo deseado mediante multiplicación por el cofactor de la curva.
\end{enumerate}

Más precisamente, estos mapeos algebraicos consisten en funciones racionales definidas sobre $\Fq$, diseñadas de modo que su imagen pertenezca a $\E(\Fq)$ con alta probabilidad o de forma determinística. Dichas funciones se construyen utilizando identidades algebraicas que garantizan la satisfacción de la ecuación de la curva.

\subsection{Firmas Schnorr sobre curvas elípticas}
\label{subsec:apendice-a-hashes-firmas-schnorr-sobre-curvas-elipticas}

Si fijamos un subgrupo cíclico de una curva elíptica $\gen{G} \subseteq \E(\Fq)$ de orden primo $r$, es natural construir sobre él primitivas criptográficas concretas. Un ejemplo es el esquema de firmas de Schnorr, que será utilizado en la tesis.

\begin{itemize}[itemsep=0.1 cm]
	\item El firmante tiene una clave privada, que es un escalar $\sk \in \Fr$, y una clave pública, que es un punto $\pk = \sk G \in \gen{G}$.
	\item Para firmar un mensaje $m \in \Strings{*}$, el firmante elige un $k \in \Fr$ aleatorio, calcula
	$R = kG$, deriva el \textit{challenge} $c = \Hash(m \concat R) \in \Fr$, define la respuesta $s = k + c \cdot \sk \in \Fr$, y publica como firma el par $\sigma = (R,s)$.
	\item La verificación consiste en recomputar $c = \Hash(m \concat R)$ y comprobar si se verifica la igualdad $sG \stackrel{?}{=} R + c\cdot \pk$. Esta igualdad se satisface si el firmante conocía la clave privada $\sk$, ya que $sG = (k + c \cdot \sk)G = kG + c \cdot \sk G = R + c \cdot \pk$.
\end{itemize}

Informalmente, producir firmas válidas sin conocer $\sk$ requiere quebrar la dificultad esperada del \ECDLP o manipular indebidamente el \textit{challenge} derivado con la función de hash. Por esa razón, Schnorr constituye una de las construcciones más limpias para entender cómo se pasa de la teoría de grupos de curvas elípticas a primitivas criptográficas concretas.
